\documentclass[uplatex,dvipdfmx,a4paper,11pt]{jsarticle}
\input{../../../templates/kaisetsu.tex}
\kaisetsuprint{熱物理学　第3章「熱力学で使う数学（偏微分・全微分・線積分）」}{2026-11-12}
%% kaisetsu-moto	problems\熱物理学\shou03\熱物理学-3.tex	kiroku\kaisetsu\熱物理学\熱物理学-3_2026-11-12_komento.tex
\begin{document}
\matome
\begin{saitenichiran}
  \hyougyou{1}{L1・計算}{◎◎◎}{S}{全問 ◎}
  \hyougyou{\hyperlink{bt-m-017-pfaff}{2}}{L1・計算}{▲◎△}{C}{(1)方針 (3)未着手}
  \hyougyou{\hyperlink{bt-m-019-joutai}{4}}{L2・計算}{◎◎◎◎△}{A}{(5)計算}
  \hyougyou{\hyperlink{bt-m-020-henkan}{5}}{L2・導出}{○○◎}{B}{(1)記述 (2)記述}
  \hyougyou{\hyperlink{bt-m-021-setsumei}{6}}{L2・説明}{△△}{C}{(1)記述 (2)記述}
  \hyougyou{\hyperlink{bt-m-022-genshi}{7}}{L2・計算}{◎◎○}{A}{(3)記述}
\end{saitenichiran}
\begin{daiji}
  \item 原始関数は $\int A\,\dd T + \int B\,\dd V$ の和ではない。一方を不定積分して「積分定数」を残りの変数の関数 $g(V)$ とおき，もう一方の係数と見比べて決める。求めたら偏微分して係数に戻るかで検算する（問題2 (1)・(3)）。
  \item 式の並びだけでは伝わらない根拠を一言書く。係数を比べてよいのは比べる微小量が独立に選べるから，線積分が $F(\mathrm{Q})-F(\mathrm{P})$ に等しいのは全微分だから（問題5 (1)(2)・問題7 (3)）。
  \item 説明の問いには，問われたことを全部答える。符号の整合は因子ごとに $\alpha>0$，$\kappa_{T}>0$ から符号を決め，積が 0 になる条件は「どちらかの因子が 0」の2通りを挙げる（問題6）。
\end{daiji}

\youtenseiri
\begin{noteshikilist}
  \item 全微分 % 出典: p.7
    \[ \dd X = \left(\pdv{X}{T}\right)_{V}\dd T + \left(\pdv{X}{V}\right)_{T}\dd V \notetag{3.1} \]
  \item 逆関数の微分 % 出典: p.8
    \[ \left(\pdv{V}{T}\right)_{p}\left(\pdv{T}{V}\right)_{p} = 1 \notetag{3.2} \]
  \item 循環関係 % 出典: p.8
    \[ \left(\pdv{p}{T}\right)_{V}\left(\pdv{T}{V}\right)_{p}\left(\pdv{V}{p}\right)_{T} = -1 \notetag{3.3} \]
  \item 熱圧力係数 % 出典: p.8
    \[ \left(\pdv{p}{T}\right)_{V} = \frac{\alpha}{\kappa_{T}} \notetag{3.4} \]
  \item 変数変換（$T$ 微分） % 出典: p.9
    \[ \left(\pdv{X}{T}\right)_{p} = \left(\pdv{X}{T}\right)_{V} + \alpha V\left(\pdv{X}{V}\right)_{T} \notetag{3.5} \]
  \item 変数変換（$p$ 微分） % 出典: p.9
    \[ \left(\pdv{X}{p}\right)_{T} = -\kappa_{T} V\left(\pdv{X}{V}\right)_{T} \notetag{3.6} \]
  \item 媒介変数表示による線積分（$s$ が $a$ から $b$ まで動くとき $C$ を描く） % 出典: p.10
    \[ \int_{C}(A\,\dd T + B\,\dd V) = \int_{a}^{b}\left(A(T(s),V(s))\dv{T}{s} + B(T(s),V(s))\dv{V}{s}\right)\dd s \notetag{3.7} \]
  \item 局所条件 % 出典: p.11
    \[ \left(\pdv{B}{T}\right)_{V} = \left(\pdv{A}{V}\right)_{T} \notetag{3.9} \]
  \item 原始関数の構成（基準点 $(T^{*},V^{*})$） % 出典: p.12
    \[ F(T,V) = \int_{T^{*}}^{T} A(T',V^{*})\,\dd T' + \int_{V^{*}}^{V} B(T,V')\,\dd V' \notetag{3.10} \]
  \item 線積分と原始関数（始点 $(T_{0},V_{0})$，終点 $(T_{N},V_{N})$） % 出典: p.12
    \[ \int_{C}(A\,\dd T + B\,\dd V) = F(T_{N},V_{N}) - F(T_{0},V_{0}) \notetag{3.11} \]
\end{noteshikilist}

\mondaihen
\level{1}
\mondaibangou{2}
\begin{mondai}{017-pfaff}{計算}{条件2による全微分の判定と，不定積分による原始関数の求め方（2変数・3変数）}
\begin{teisuu}
  気体定数 & $R$ & \qty{8.314}{\joule\per\kelvin\per\mole}\\
\end{teisuu}
$b$，$c$ を定数とする。
  \begin{shoumon}
    \item 次の Pfaff 形式が全微分かどうかを判定し，全微分なら原始関数 $F(T,V)$ を求めよ。
    \[ \omega_{1} = (2TV^{2} + b)\,\dd T + \left(2T^{2}V - \frac{c}{V^{2}}\right)\dd V \]
    \item $\omega_{2} = V\,\dd T - T\,\dd V$ について同じことをせよ。
    \item 3変数 $(T,V,n)$ の Pfaff 形式
    \[ \omega_{3} = nV\,\dd T + nT\,\dd V + (TV + 2n)\,\dd n \]
    について，一般形の条件2がすべて成り立つことを確かめ，原始関数 $F(T,V,n)$ を求めよ。
  \end{shoumon}

\begin{kotaeran}
  \machigai{1}{▲}{方針} 条件2で全微分と判定したところまでは正しい。原始関数を $F=\int(2TV^{2}+b)\,\dd T+\int(2T^{2}V-c/V^{2})\,\dd V=2T^{2}V^{2}+bT+c/V$ とした。2つの係数をそれぞれ積分して足すと，両方に含まれる $T^{2}V^{2}$ を2回数えてしまう。実際，この $F$ を $T$ で微分すると $4TV^{2}+b$ になり，$\dd T$ の係数 $2TV^{2}+b$ に戻らない。正しくは $T$ で不定積分して $F=T^{2}V^{2}+bT+g(V)$ とおき，$V$ で微分して $B$ と比べると $g'(V)=-c/V^{2}$ なので $F=T^{2}V^{2}+bT+c/V+(\text{定数})$。
  \machigai{3}{△}{未着手} $(T,V)$，$(T,n)$，$(V,n)$ の3組で条件2を確かめたところは正しいが，原始関数 $F(T,V,n)$ を求めていない。(1) と同じ手順で $F=nTV+h(V,n)$ から始め，$V$ 微分を $nT$ と，$n$ 微分を $TV+2n$ と順に見比べると $F=nTV+n^{2}+(\text{定数})$。
\end{kotaeran}
\begin{rikai}
  \item $(\partial F/\partial T)_{V}=A$ を $T$ で積分するとき $V$ は固定していたので，「積分定数」は $V$ の任意関数 $g(V)$ になる。$g$ は $(\partial F/\partial V)_{T}=B$ と見比べて決める。$B$ を $V$ で積分した結果をそのまま足すと，$A$ からすでに出ている項が重なる。
  \item もう1つの作り方は，基準点から $V=V^{*}$ のまま $T$ 方向，続けて $T$ を固定して $V$ 方向に進む線積分 $F=\int_{T^{*}}^{T}A(T',V^{*})\,\dd T'+\int_{V^{*}}^{V}B(T,V')\,\dd V'$。問題7 (2) ではこれができている。
  \item 求めた $F$ は必ず偏微分して，$\dd T$・$\dd V$（・$\dd n$）の係数に戻ることを確かめる。
\end{rikai}
\end{mondai}

\level{2}
\mondaibangou{4}
\begin{mondai}{019-joutai}{計算}{状態方程式から得る体膨張率・等温圧縮率と，逆関数の微分・循環関係・熱圧力係数の関係式}
\begin{teisuu}
  気体定数 & $R$ & \qty{8.314}{\joule\per\kelvin\per\mole}\\
\end{teisuu}
物質量 $n$ を固定した気体が，状態方程式 $p(V - nb) = nRT$（$b$ は正の定数，$V > nb$）に従うとする。
  \begin{shoumon}
    \item 体膨張率 $\alpha$ と等温圧縮率 $\kappa_{T}$ を求め，$R$ を含まない形（$T$，$p$，$V$，$n$，$b$ で表す形）にせよ。
    \item $(\partial p/\partial T)_{V}$ を状態方程式から直接求め，\noteshiki{3.4} が成り立つことを確かめよ。
    \item $(\partial p/\partial V)_{T}$ を状態方程式から直接求め，$-1/(\kappa_{T}V)$ に等しいことを確かめよ。
    \item $(\partial p/\partial T)_{V}$，$(\partial T/\partial V)_{p}$，$(\partial V/\partial p)_{T}$ をそれぞれ直接求め，\noteshiki{3.3} を確かめよ。
    \item $nb/V = 0.0200$，$T = \qty{300}{\kelvin}$ の状態での $\alpha$ を有効数字3桁で求め，同じ温度の理想気体の値と比べよ。
  \end{shoumon}

\begin{kotaeran}
  \machigai{5}{△}{計算} $\alpha=(1/T)(1-nb/V)$ までは正しいが，値を $2.67\times10^{-4}$ とした。正しくは $(1-0.0200)/\qty{300}{\kelvin}=\qty{3.27e-3}{\per\kelvin}$。$2.67\times10^{-4}$ は $0.08/300$ に当たるので，$1-0.0200=0.980$ の引き算を誤ったとみられる。単位も無く，問われている同じ温度の理想気体の値 $1/T=\qty{3.33e-3}{\per\kelvin}$ との比較（この気体は $0.980$ 倍で 2.0\,\% 小さい）も書かれていない。
\end{kotaeran}
\begin{rikai}
  \item $nb/V=0.0200$ は小さな補正なので，$\alpha$ は理想気体の $1/T$ に近いはずである。$b\to0$ で理想気体に戻ることを使って，答えが $1/T$ と同じ桁かを確かめる。
  \item 数値の答えには単位（\unit{\per\kelvin}）を付ける。「比べよ」には比（$0.980$ 倍）か差で答える。
\end{rikai}
\end{mondai}

\mondaibangou{5}
\begin{mondai}{020-henkan}{導出}{全微分の代入による偏微分係数の変数変換}
\begin{teisuu}
  気体定数 & $R$ & \qty{8.314}{\joule\per\kelvin\per\mole}\\
\end{teisuu}
物質量 $n$ を固定した物質の状態量 $X$ を考える。
  \begin{shoumon}
    \item $\dd V = \alpha V\,\dd T - \kappa_{T} V\,\dd p$ を示し，これを $X(T,V)$ の全微分に代入して \noteshiki{3.5}・\noteshiki{3.6} を導け。
    \item 独立変数を $(V,p)$ に選んだときの $(\partial X/\partial V)_{p}$ と $(\partial X/\partial p)_{V}$ を，$\alpha$，$\kappa_{T}$，$V$ と $(\partial X/\partial T)_{V}$，$(\partial X/\partial V)_{T}$ で表せ（$\alpha \neq 0$ とする）。さらに $X = T$ とおいた結果が \noteshiki{3.2}・\noteshiki{3.4} と整合することを確かめよ。
    \item 状態方程式 $p(V - nb) = nRT$ に従う気体（$b$ は正の定数。体膨張率は $\alpha = (V - nb)/(TV)$）で，$X(T,V) = cT + k/V$（$c$，$k$ は定数）とする。$(\partial X/\partial T)_{p}$ を \noteshiki{3.5} で求め，$V$ を $(T,p)$ で表して直接微分した結果と一致することを確かめよ。
  \end{shoumon}

\begin{kotaeran}
  \machigai{1}{○}{記述} $\dd V=\alpha V\,\dd T-\kappa_{T}V\,\dd p$ の導出と代入，得た $(\partial X/\partial T)_{p}=(\partial X/\partial T)_{V}+\alpha V(\partial X/\partial V)_{T}$，$(\partial X/\partial p)_{T}=-\kappa_{T}V(\partial X/\partial V)_{T}$ は正しい。代入した $\dd X$ からすぐ $(\partial X/\partial T)_{p}=\cdots$ と書いていて，$X$ を $(T,p)$ の関数とみた全微分 $\dd X=(\partial X/\partial T)_{p}\dd T+(\partial X/\partial p)_{T}\dd p$ と，係数を比べてよい理由（$\dd T$ と $\dd p$ が独立に選べる）が無い。
  \machigai{2}{○}{記述} $(V,p)$ の全微分に $\dd V$ を代入し，$\dd T$・$\dd p$ の係数を (1) の結果と比べて $(\partial X/\partial V)_{p}$，$(\partial X/\partial p)_{V}$ を出す方法は，解答と違うが正しい。$X=T$ での確かめも正しい。(1) と同じく係数比較の根拠が無く，$1/(\alpha V)$ を作るときに $\alpha\neq0$ を使ったことも書かれていない。
\end{kotaeran}
\begin{rikai}
  \item $\dd T$，$\dd V$，$\dd p$ は状態方程式で結ばれていて，独立に選べるのは2つだけである。係数を比べるときは「どの2つが独立か」を一言書く。3つを同時に独立とみると誤る。
  \item 割り算をしたら，割った量が 0 でないこと（ここでは $\alpha\neq0$）を書く。水の \qty{4}{\degreeCelsius} 付近のように $\alpha=0$ となる状態では，$(V,p)$ を独立変数に選べない。
\end{rikai}
\end{mondai}

\mondaibangou{6}
\begin{mondai}{021-setsumei}{説明}{何を固定した偏微分かの区別（循環関係の符号と，定圧・定積での温度微分の違い）}
\begin{teisuu}
  気体定数 & $R$ & \qty{8.314}{\joule\per\kelvin\per\mole}\\
\end{teisuu}
\begin{shoumon}
    \item ある学生は「$(\partial p/\partial T)_{V}(\partial T/\partial V)_{p}(\partial V/\partial p)_{T}$ は $\partial T$，$\partial V$，$\partial p$ を約分すれば $1$ になる」と考えた。この考えの誤りを説明せよ。さらに，3つの因子をそれぞれ $\alpha$，$\kappa_{T}$，$V$ で表し，正しい値 $-1$ の符号が，通常の物質（$\alpha > 0$，$\kappa_{T} > 0$）の性質と整合していることを説明せよ。
    \item 同じ状態で $(\partial X/\partial T)_{p}$ と $(\partial X/\partial T)_{V}$ が等しくなるのはどのような場合か。\noteshiki{3.5} を使って条件を挙げ，それぞれの意味を $(T,V)$ 平面上の等積線と等圧線の言葉で説明せよ。
  \end{shoumon}

\begin{kotaeran}
  \machigai{1}{△}{記述} 3つの因子を $\alpha/\kappa_{T}$，$1/(\alpha V)$，$-\kappa_{T}V$ と表して積が $-1$ になることは正しい。約分できない理由は「固定した変数について考えていない」だけで，3つが等積・等圧・等温という別々の変化に沿った比なので，分母と分子の $\partial T$ などが同じ量ではない，まで書かれていない。また，問われている「$-1$ の符号が $\alpha>0$，$\kappa_{T}>0$ と整合すること」の説明が無い。積 $(\alpha/\kappa_{T})\cdot(1/(\alpha V))\cdot(-\kappa_{T}V)$ は $\alpha$，$\kappa_{T}$ の符号によらず $-1$ になるので，この計算だけでは整合の説明にならない。
  \machigai{2}{△}{記述} $(\partial X/\partial T)_{p}=(\partial X/\partial T)_{V}+\alpha V(\partial X/\partial V)_{T}$ から $(\partial X/\partial V)_{T}=0$ を挙げたのは正しいが，差 $\alpha V(\partial X/\partial V)_{T}$ が 0 になるもう1つの場合 $\alpha=0$ が式で挙がっていない（図に添えた「等積線と等圧線の傾きが同じ」は $\alpha=0$ の意味に当たるが，$\alpha$ と結び付けていない）。$(\partial X/\partial V)_{T}=0$ の意味も「$X$ が $V$ によらない」にとどまり，等積線・等圧線の言葉での説明が無い。
\end{kotaeran}
\begin{rikai}
  \item 符号は因子ごとに決める。$\alpha>0$（定圧で温めると膨張）なので $(\partial T/\partial V)_{p}=1/(\alpha V)>0$，$\kappa_{T}>0$（等温で圧縮すると体積が減る）なので $(\partial V/\partial p)_{T}=-\kappa_{T}V<0$，その結果 $(\partial p/\partial T)_{V}=\alpha/\kappa_{T}>0$。正・正・負で積は負になり，$-1$ の符号と合う。
  \item 積が 0 になるのは，どちらかの因子が 0 のとき。$V>0$ なので $\alpha V(\partial X/\partial V)_{T}=0$ は「$\alpha=0$ または $(\partial X/\partial V)_{T}=0$」の2通り。
  \item $(T,V)$ 平面（横軸 $T$，縦軸 $V$）で，等積線は水平線，等圧線の傾きは $(\partial V/\partial T)_{p}=\alpha V$。$\alpha=0$ なら等圧線がその点で等積線に接し，定圧で温めても体積が1次で変わらない。$(\partial X/\partial V)_{T}=0$ なら，等圧線が等積線からずれて体積が $\alpha V\,\dd T$ 変わっても，$X$ がその方向に変わらないので差が出ない。
\end{rikai}
\end{mondai}

\mondaibangou{7}
\begin{mondai}{022-genshi}{計算}{線積分による原始関数の構成と，線積分が原始関数の差に等しいこと}
\begin{teisuu}
  気体定数 & $R$ & \qty{8.314}{\joule\per\kelvin\per\mole}\\
\end{teisuu}
Pfaff 形式 $\omega = V^{2}\,\dd T + 2TV\,\dd V$ を考える。$T_{0}$，$V_{0}$ は正の定数とする。
  \begin{shoumon}
    \item $\omega$ が条件2を満たすことを確かめよ。
    \item 基準点を $(T^{*},V^{*})$ として，\noteshiki{3.10} の構成で原始関数 $F(T,V)$ を求めよ。さらに $(\partial F/\partial T)_{V} = V^{2}$ と $(\partial F/\partial V)_{T} = 2TV$ を確かめよ。
    \item 点 $\mathrm{P} = (T_{0},V_{0})$ から点 $\mathrm{Q} = (2T_{0},V_{0}/2)$ に至る曲線 $C$：$(T,V) = (T_{0}(1+s),V_{0}(1-s/2))$（$0 \le s \le 1$）に沿った $\omega$ の線積分を，原始関数を使わずに媒介変数表示で直接計算せよ。その値が $F(\mathrm{Q}) - F(\mathrm{P})$ に等しいことを確かめよ。
  \end{shoumon}

\begin{kotaeran}
  \machigai{3}{○}{記述} 媒介変数での計算，値 $-T_{0}V_{0}^{2}/2$，$F(\mathrm{Q})-F(\mathrm{P})$ との一致は正しい。書き出しを $F(T,V)=\int_{C}\omega$ としたが，線積分は曲線 $C$ で決まる1つの数で，$(T,V)$ の関数 $F$ ではない（等しいのは $F(\mathrm{Q})-F(\mathrm{P})$）。また，両者が等しくなる理由（(1) で条件2を確かめた全微分なので，線積分は原始関数の終点と始点の値の差になる。基準点による定数 $T^{*}V^{*2}$ は差で消える）が書かれていない。
\end{kotaeran}
\begin{rikai}
  \item 線積分 $\int_{C}\omega$ が $F(\text{終点})-F(\text{始点})$ に等しいのは，$\omega$ が全微分（原始関数をもつ）ときだけである。名前を付けるなら $I_{C}$ などとし，原始関数 $F$ と区別する。
  \item (2) で最初に第2項の $T$ を $T^{*}$ にしかけて自分で直せている。原始関数の構成の第2項 $\int_{V^{*}}^{V}B(T,V')\,\dd V'$ の $T$ は終点の $T$ である。
\end{rikai}
\end{mondai}

\kaitouhen
\level{1}
\kaisetsubangou{2}
\begin{kaisetsu}{017-pfaff}
\houshin 条件2（\noteshiki{3.9}，3変数では一般形）で判定し，全微分なら $(\partial F/\partial T)_{V} = A$ を $T$ で不定積分して，積分定数を残りの変数の関数とおく。それをもう一方の条件と見比べて決める。
\keisan \shoumonkai{1} $A = 2TV^{2} + b$，$B = 2T^{2}V - c/V^{2}$ とおくと
\[ \left(\pdv{A}{V}\right)_{T} = 4TV, \qquad \left(\pdv{B}{T}\right)_{V} = 4TV \]
で恒等的に等しいので，$\omega_{1}$ は全微分である。$(\partial F/\partial T)_{V} = 2TV^{2} + b$ を $T$ で不定積分すると，$g$ を $V$ だけの関数として
\[ F = T^{2}V^{2} + bT + g(V) \]
これを $V$ で微分して $B$ と比べる。
\[ \left(\pdv{F}{V}\right)_{T} = 2T^{2}V + g'(V) = 2T^{2}V - \frac{c}{V^{2}} \]
よって $g'(V) = -c/V^{2}$，$g(V) = c/V + (\text{定数})$。
\shoumonkai{2} $A = V$，$B = -T$ なので
\[ \left(\pdv{A}{V}\right)_{T} = 1, \qquad \left(\pdv{B}{T}\right)_{V} = -1 \]
どの点でも等しくないので条件2が成り立たず，$\omega_{2}$ は全微分でない（不完全微分。原始関数を持たない）。
\shoumonkai{3} $X_{1} = T$，$X_{2} = V$，$X_{3} = n$ とし，$A_{1} = nV$，$A_{2} = nT$，$A_{3} = TV + 2n$ とおく。3つの組を調べる（偏微分は残りの2変数を固定したもの）。
\begin{align*}
&(T,V)\text{ の組：} \pdv{A_{1}}{V} = n, \quad \pdv{A_{2}}{T} = n,\\
&(T,n)\text{ の組：} \pdv{A_{1}}{n} = V, \quad \pdv{A_{3}}{T} = V,\\
&(V,n)\text{ の組：} \pdv{A_{2}}{n} = T, \quad \pdv{A_{3}}{V} = T
\end{align*}
すべて等しいので全微分である。$\partial F/\partial T = nV$ を $T$ で不定積分すると，$h$ を $(V,n)$ の関数として $F = nTV + h(V,n)$。$V$ で微分して $A_{2}$ と比べると
\[ \pdv{F}{V} = nT + \pdv{h}{V} = nT \]
より $h$ は $n$ だけの関数 $h(n)$。$n$ で微分して $A_{3}$ と比べると
\[ \pdv{F}{n} = TV + h'(n) = TV + 2n \]
より $h(n) = n^{2} + (\text{定数})$。
\kotae{\shoumonkai{1} 全微分。$F = T^{2}V^{2} + bT + c/V + (\text{定数})$ \shoumonkai{2} 不完全微分（$(\partial A/\partial V)_{T} = 1 \neq -1 = (\partial B/\partial T)_{V}$） \shoumonkai{3} $F = nTV + n^{2} + (\text{定数})$}
\begin{ronten}
  \item \shoumonkai{1} 条件2が恒等的に成り立つことから，三条件の同値性により全微分（原始関数をもつ）と判定したことを書く。
  \item \shoumonkai{1} 不定積分の積分定数を，定数ではなく固定していた変数 $V$ の関数 $g(V)$ とおく理由を述べる。
  \item \shoumonkai{2} 条件2が成り立たないことから，三条件の同値性により原始関数をもたないと結論する。
  \item \shoumonkai{3} 3変数では $(T,V)$，$(T,n)$，$(V,n)$ の3組すべてで条件2を確かめる必要があることを述べる。
  \item \shoumonkai{3} 積分定数が残りの2変数の関数 $h(V,n)$ になることと，$A_{2}$ と比べて $h$ が $V$ によらないと分かってから $A_{3}$ と比べることを書く。
\end{ronten}
\kenzan 求めた $F$ を微分して戻す。(1)：$\dd F = (2TV^{2} + b)\,\dd T + (2T^{2}V - c/V^{2})\,\dd V$ で $\omega_{1}$ に一致する。(3)：$\dd F = nV\,\dd T + nT\,\dd V + (TV + 2n)\,\dd n$ で $\omega_{3}$ に一致する。
\tsumazuki 不定積分の「積分定数」は，固定していた変数の任意関数である。定数とおくと，もう一方の条件と見比べたときに項が合わなくなる。
\end{kaisetsu}

\level{2}
\kaisetsubangou{4}
\begin{kaisetsu}{019-joutai}
\houshin 状態方程式を $V$ について解いて $\alpha$，$\kappa_{T}$ を，$p$ について解いて $(\partial p/\partial T)_{V}$，$(\partial p/\partial V)_{T}$ を，$T$ について解いて $(\partial T/\partial V)_{p}$ を求める。$R$ は状態方程式で消す。
\keisan \shoumonkai{1} $V = nRT/p + nb$ より
\[ \left(\pdv{V}{T}\right)_{p} = \frac{nR}{p}, \qquad \left(\pdv{V}{p}\right)_{T} = -\frac{nRT}{p^{2}} \]
状態方程式から $nR/p = (V - nb)/T$，$nRT/p = V - nb$ なので
\[ \alpha = \frac{1}{V}\cdot\frac{nR}{p} = \frac{V - nb}{TV}, \qquad \kappa_{T} = \frac{1}{V}\cdot\frac{nRT}{p^{2}} = \frac{V - nb}{pV} \]
\shoumonkai{2} $p = nRT/(V - nb)$ を $V$ 固定で $T$ 微分すると
\[ \left(\pdv{p}{T}\right)_{V} = \frac{nR}{V - nb} = \frac{p}{T} \]
一方，(1) から
\[ \frac{\alpha}{\kappa_{T}} = \frac{V - nb}{TV}\cdot\frac{pV}{V - nb} = \frac{p}{T} \]
で一致し，\noteshiki{3.4} が成り立つ。
\shoumonkai{3} $T$ 固定で $V$ 微分すると
\[ \left(\pdv{p}{V}\right)_{T} = -\frac{nRT}{(V - nb)^{2}} = -\frac{p}{V - nb} \]
一方，(1) から
\[ -\frac{1}{\kappa_{T}V} = -\frac{pV}{(V - nb)V} = -\frac{p}{V - nb} \]
で一致する（\noteshiki{3.2} を $T$ 固定の $p$ と $V$ に使ったもの）。
\shoumonkai{4} (2) より $(\partial p/\partial T)_{V} = p/T$。$T = p(V - nb)/(nR)$ を $p$ 固定で $V$ 微分し，$p/(nR) = T/(V - nb)$ を使うと
\[ \left(\pdv{T}{V}\right)_{p} = \frac{p}{nR} = \frac{T}{V - nb} \]
(1) より $(\partial V/\partial p)_{T} = -nRT/p^{2} = -(V - nb)/p$。3つを掛けると
\[ \frac{p}{T}\cdot\frac{T}{V - nb}\cdot\left(-\frac{V - nb}{p}\right) = -1 \]
で \noteshiki{3.3} が成り立つ。
\shoumonkai{5} (1) より $\alpha = (1 - nb/V)/T$。数値を入れると
\[ \alpha = \frac{1 - 0.0200}{\qty{300}{\kelvin}} = \qty{3.27e-3}{\per\kelvin} \]
同じ温度の理想気体では $\alpha = 1/T = \qty{3.33e-3}{\per\kelvin}$ で，この気体の値はその $0.980$ 倍（2.0\,\% 小さい）。
\kotae{\shoumonkai{1} $\alpha = (V - nb)/(TV)$，$\kappa_{T} = (V - nb)/(pV)$ \shoumonkai{2} $(\partial p/\partial T)_{V} = p/T = \alpha/\kappa_{T}$ \shoumonkai{3} $(\partial p/\partial V)_{T} = -p/(V - nb) = -1/(\kappa_{T}V)$ \shoumonkai{4} $p/T$，$T/(V - nb)$，$-(V - nb)/p$ で積は $-1$ \shoumonkai{5} $\alpha = \qty{3.27e-3}{\per\kelvin}$（理想気体の \qty{3.33e-3}{\per\kelvin} より 2.0\,\% 小さい）}
\kenzan 次元：$\alpha$ は $\unit{\per\kelvin}$，$\kappa_{T}$ は $\unit{\per\pascal}$ で定義と合う。極限 $b \to 0$ で $\alpha \to 1/T$，$\kappa_{T} \to 1/p$ となり，理想気体の結果に戻る。
\tsumazuki \noteshiki{3.3} の左辺を約分すると $+1$ になり，符号を誤る。各因子は固定する変数が違う別々の比である。
\end{kaisetsu}

\kaisetsubangou{5}
\begin{kaisetsu}{020-henkan}
\houshin 変数を切り替えるには，一方の組の全微分にもう一方の組の全微分を代入し，独立に動かせる微小量の係数を比べる。$\dd T$ と $\dd p$（または $\dd V$ と $\dd p$）は独立に選べるので，係数が一致しなければならない。
\keisan \shoumonkai{1} $V$ を $(T,p)$ の関数とみた全微分に，$\alpha$，$\kappa_{T}$ の定義を入れる。
\[ \dd V = \left(\pdv{V}{T}\right)_{p}\dd T + \left(\pdv{V}{p}\right)_{T}\dd p = \alpha V\,\dd T - \kappa_{T}V\,\dd p \]
これを \noteshiki{3.1} の $\dd V$ に代入すると
\[ \dd X = \left[\left(\pdv{X}{T}\right)_{V} + \alpha V\left(\pdv{X}{V}\right)_{T}\right]\dd T - \kappa_{T}V\left(\pdv{X}{V}\right)_{T}\dd p \]
一方，$X$ を $(T,p)$ の関数とみると
\[ \dd X = \left(\pdv{X}{T}\right)_{p}\dd T + \left(\pdv{X}{p}\right)_{T}\dd p \]
$\dd p = 0$ とおけば $\dd T$ の係数，$\dd T = 0$ とおけば $\dd p$ の係数が等しくなるので，\noteshiki{3.5}・\noteshiki{3.6} が得られる。
\shoumonkai{2} $\dd V = \alpha V\,\dd T - \kappa_{T}V\,\dd p$ を $\dd T$ について解く（$\alpha \neq 0$）。
\[ \dd T = \frac{1}{\alpha V}\dd V + \frac{\kappa_{T}}{\alpha}\dd p \]
これを \noteshiki{3.1} の $\dd T$ に代入すると
\[ \dd X = \left[\left(\pdv{X}{V}\right)_{T} + \frac{1}{\alpha V}\left(\pdv{X}{T}\right)_{V}\right]\dd V + \frac{\kappa_{T}}{\alpha}\left(\pdv{X}{T}\right)_{V}\dd p \]
$X$ を $(V,p)$ の関数とみた全微分と係数を比べて
\[ \left(\pdv{X}{V}\right)_{p} = \left(\pdv{X}{V}\right)_{T} + \frac{1}{\alpha V}\left(\pdv{X}{T}\right)_{V}, \qquad \left(\pdv{X}{p}\right)_{V} = \frac{\kappa_{T}}{\alpha}\left(\pdv{X}{T}\right)_{V} \]
$X = T$ とおくと $(\partial T/\partial T)_{V} = 1$，$(\partial T/\partial V)_{T} = 0$ なので
\[ \left(\pdv{T}{V}\right)_{p} = \frac{1}{\alpha V}, \qquad \left(\pdv{T}{p}\right)_{V} = \frac{\kappa_{T}}{\alpha} \]
第1式は $(\partial V/\partial T)_{p} = \alpha V$ の逆数で \noteshiki{3.2} と，第2式は \noteshiki{3.4} の逆数で整合する。
\shoumonkai{3} $(\partial X/\partial T)_{V} = c$，$(\partial X/\partial V)_{T} = -k/V^{2}$ を \noteshiki{3.5} に入れる。
\[ \left(\pdv{X}{T}\right)_{p} = c + \frac{V - nb}{TV}\cdot V\cdot\left(-\frac{k}{V^{2}}\right) = c - \frac{k(V - nb)}{TV^{2}} \]
直接計算では，$V = nRT/p + nb$ を代入した $X = cT + k/V(T,p)$ を $p$ 固定で微分する。$(\partial V/\partial T)_{p} = nR/p$ なので
\[ \left(\pdv{X}{T}\right)_{p} = c - \frac{k}{V^{2}}\cdot\frac{nR}{p} \]
状態方程式から $nR/p = (V - nb)/T$ なので，上と一致する。
\kotae{\shoumonkai{1} 本文の導出のとおり \shoumonkai{2} $(\partial X/\partial V)_{p} = (\partial X/\partial V)_{T} + (\partial X/\partial T)_{V}/(\alpha V)$，$(\partial X/\partial p)_{V} = (\kappa_{T}/\alpha)(\partial X/\partial T)_{V}$。$X = T$ で $(\partial T/\partial V)_{p} = 1/(\alpha V)$，$(\partial T/\partial p)_{V} = \kappa_{T}/\alpha$ \shoumonkai{3} $(\partial X/\partial T)_{p} = c - k(V - nb)/(TV^{2})$}
\begin{ronten}
  \item \shoumonkai{1} $V$ を $(T,p)$ の関数とみて全微分を書き，その係数が $\alpha$，$\kappa_{T}$ の定義から $\alpha V$，$-\kappa_{T}V$ になることを述べる。
  \item 係数を比べてよいのは，比べる2つの微小量（(1) では $\dd T$ と $\dd p$，(2) では $\dd V$ と $\dd p$）が独立に選べるからであることを書く。
  \item \shoumonkai{2} $\dd T$ について解くときに $\alpha \neq 0$ を使ったことを明記する。
\end{ronten}
\kenzan (2) で $X = V$ とおくと $(\partial V/\partial V)_{p} = 1$，$(\partial V/\partial p)_{V} = 0$ となり正しい。(3) は $k = 0$ で $c$ に戻り，$b \to 0$ で $c - k/(TV)$ となる（理想気体では $\alpha V = V/T$）。
\tsumazuki 係数を比べてよいのは，比べる2つの微小量が独立に選べるときだけである。$\dd T$，$\dd V$，$\dd p$ の3つは状態方程式で結ばれているので，3つを同時に独立とみてはいけない。
\end{kaisetsu}

\kaisetsubangou{6}
\begin{kaisetsu}{021-setsumei}
\houshin 偏微分係数は「どの変数を固定した変化に沿った比か」で決まる。各因子を $\alpha$，$\kappa_{T}$，$V$ で表して符号を読み，(2) は \noteshiki{3.5} の右辺第2項を調べる。
\keisan \shoumonkai{1} $(\partial p/\partial T)_{V}$ は等積変化，$(\partial T/\partial V)_{p}$ は等圧変化，$(\partial V/\partial p)_{T}$ は等温変化に沿った無限小変化の比である。第1因子の分母の $\partial T$ と第2因子の分子の $\partial T$ は，違う過程に沿った温度変化なので同じ量ではなく，約分できない。各因子は，\noteshiki{3.4}，定義と \noteshiki{3.2}，定義から
\[ \left(\pdv{p}{T}\right)_{V} = \frac{\alpha}{\kappa_{T}}, \qquad \left(\pdv{T}{V}\right)_{p} = \frac{1}{\alpha V}, \qquad \left(\pdv{V}{p}\right)_{T} = -\kappa_{T}V \]
となり，積は
\[ \frac{\alpha}{\kappa_{T}}\cdot\frac{1}{\alpha V}\cdot(-\kappa_{T}V) = -1 \]
通常の物質では，定圧で温めると膨張する（$\alpha > 0$）ので第2因子は正，等温で圧縮すると体積が減る（$\kappa_{T} > 0$）ので第3因子は負，その結果，体積を一定に保って温めると圧力が上がり第1因子は正になる。正・正・負の積なので負であり，$-1$ の符号と整合する。
\shoumonkai{2} \noteshiki{3.5} から
\[ \left(\pdv{X}{T}\right)_{p} - \left(\pdv{X}{T}\right)_{V} = \alpha V\left(\pdv{X}{V}\right)_{T} \]
両者が等しいのは，その状態で $\alpha = 0$ か $(\partial X/\partial V)_{T} = 0$ のときである。$(T,V)$ 平面で，等積線は $V$ 一定の直線，等圧線は傾き $(\partial V/\partial T)_{p} = \alpha V$ の曲線である。定圧で温度を $\dd T$ 上げると，等積の変化に加えて体積が $\alpha V\,\dd T$ だけ変わり，そのぶん $X$ が $(\partial X/\partial V)_{T}\,\alpha V\,\dd T$ だけ余分に変わる。$\alpha = 0$ なら等圧線がその点で等積線に接し，2つの変化が1次の範囲で同じになる。$(\partial X/\partial V)_{T} = 0$ なら等圧線と等積線がずれても $X$ が $V$ 方向に変わらないので，差が出ない（例：理想気体で $X$ が $T$ だけの関数のとき）。Level 1 の最初の問題（理想気体で $X = T^{2}/V$）では $\alpha = 1/T \neq 0$，$(\partial X/\partial V)_{T} = -T^{2}/V^{2} \neq 0$ なので，2つの微分が違った（$T/V$ と $2T/V$）。
\kotae{\shoumonkai{1} 3つの因子は等積・等圧・等温という別々の変化に沿った比なので約分できない。$(\alpha/\kappa_{T})\cdot(1/(\alpha V))\cdot(-\kappa_{T}V) = -1$ で，$\alpha > 0$，$\kappa_{T} > 0$ から符号は正・正・負となり整合する \shoumonkai{2} $\alpha = 0$（等圧線が等積線に接する）か $(\partial X/\partial V)_{T} = 0$（$X$ が体積方向に変わらない）のとき}
\begin{ronten}
  \item \shoumonkai{1} 3つの因子が等積・等圧・等温という別々の変化に沿った比で，分母と分子に現れる同じ文字の微小量が同じ量ではないので約分できないことを述べる。
  \item \shoumonkai{1} 各因子を $\alpha$，$\kappa_{T}$，$V$ で表すときに使った関係式を示し，$\alpha > 0$，$\kappa_{T} > 0$，$V > 0$ から因子の符号を1つずつ決めて積が負になることを書く。
  \item \shoumonkai{2} \noteshiki{3.5} の差が $\alpha V$ と $(\partial X/\partial V)_{T}$ の積で $V > 0$ なので，条件が $\alpha = 0$ か $(\partial X/\partial V)_{T} = 0$ に限られることを述べる。
  \item \shoumonkai{2} 等圧線の傾きが $\alpha V$ であることを使って，それぞれの条件を等積線と等圧線の関係で説明する。
\end{ronten}
\kenzan 理想気体（$\alpha = 1/T$，$\kappa_{T} = 1/p$）で (1) の3因子を計算すると $p/T$，$T/V$，$-V/p$ で，積は $-1$ である。
\tsumazuki 水は約 \qty{4}{\degreeCelsius}（1気圧）で密度が最大になり，そこでは $\alpha = 0$ である。「通常の物質」で $\alpha > 0$ としたのはこのような例外を除くためである。
\end{kaisetsu}

\kaisetsubangou{7}
\begin{kaisetsu}{022-genshi}
\houshin (2) は \noteshiki{3.10} の2本の線分に沿って積分する。(3) は媒介変数表示で $s$ の積分にし（\noteshiki{3.7}），\noteshiki{3.11} と比べる。
\keisan \shoumonkai{1} $A = V^{2}$，$B = 2TV$ より
\[ \left(\pdv{A}{V}\right)_{T} = 2V, \qquad \left(\pdv{B}{T}\right)_{V} = 2V \]
で恒等的に等しい。
\shoumonkai{2} \noteshiki{3.10} より
\[ F(T,V) = \int_{T^{*}}^{T} V^{*2}\,\dd T' + \int_{V^{*}}^{V} 2TV'\,\dd V' = V^{*2}(T - T^{*}) + T(V^{2} - V^{*2}) = TV^{2} - T^{*}V^{*2} \]
これを微分すると $(\partial F/\partial T)_{V} = V^{2}$，$(\partial F/\partial V)_{T} = 2TV$ で，$A$，$B$ に一致する。
\shoumonkai{3} $T = T_{0}(1+s)$，$V = V_{0}(1 - s/2)$ より $\dd T/\dd s = T_{0}$，$\dd V/\dd s = -V_{0}/2$。被積分関数は
\begin{align*}
A\dv{T}{s} + B\dv{V}{s} &= V_{0}^{2}\left(1 - \frac{s}{2}\right)^{2}T_{0} + 2T_{0}(1+s)V_{0}\left(1 - \frac{s}{2}\right)\left(-\frac{V_{0}}{2}\right)\\
&= T_{0}V_{0}^{2}\left(1 - \frac{s}{2}\right)\left[\left(1 - \frac{s}{2}\right) - (1+s)\right]\\
&= -\frac{3}{2}T_{0}V_{0}^{2}\,s\left(1 - \frac{s}{2}\right)
\end{align*}
これを $s$ について $0$ から $1$ まで積分する。
\[ \int_{C}\omega = -\frac{3}{2}T_{0}V_{0}^{2}\left[\frac{s^{2}}{2} - \frac{s^{3}}{6}\right]_{0}^{1} = -\frac{3}{2}T_{0}V_{0}^{2}\cdot\frac{1}{3} = -\frac{1}{2}T_{0}V_{0}^{2} \]
一方，(2) の $F$ では
\[ F(\mathrm{Q}) - F(\mathrm{P}) = 2T_{0}\cdot\frac{V_{0}^{2}}{4} - T_{0}V_{0}^{2} = -\frac{1}{2}T_{0}V_{0}^{2} \]
（定数 $T^{*}V^{*2}$ は差で消える）で一致する。
\kotae{\shoumonkai{1} $(\partial A/\partial V)_{T} = (\partial B/\partial T)_{V} = 2V$ \shoumonkai{2} $F = TV^{2} - T^{*}V^{*2}$ \shoumonkai{3} $C$ に沿った $\omega$ の線積分は $-T_{0}V_{0}^{2}/2 = F(\mathrm{Q}) - F(\mathrm{P})$}
\begin{ronten}
  \item \shoumonkai{2} \noteshiki{3.10} の第2項では，$T$ を終点の値に固定したまま $V'$ について積分することを明記する。
  \item \shoumonkai{3} 線積分が $F(\mathrm{Q}) - F(\mathrm{P})$ に等しくなるのは (1) で条件2を確かめた全微分だからであり，基準点による定数は差で消えることを述べる。
\end{ronten}
\kenzan 次元はどれも $[T][V]^{2}$ でそろう。基準点を変えると $F$ は定数だけ変わり，差 $F(\mathrm{Q}) - F(\mathrm{P})$ は変わらない。$s = 0$，$1$ で被積分関数が $0$，負になり，積分が負になるのとも合う。
\tsumazuki \noteshiki{3.10} の第2項では，$T$ は終点の $T$（固定した値）で，$T^{*}$ ではない。ここを $T^{*}$ にすると $F$ が $V$ の項で合わなくなる。
\end{kaisetsu}

\end{document}
